


Conforme destaca \citeonline{aniche_2014}, teste é o processo de execução de um produto com o objetivo de verificar se ele atende às especificações definidas e funciona corretamente no ambiente para o qual foi projetado. Além disso, testes é uma etapa essencial para o desenvolvimento de software, atualmente os Estados Unidos estimam que bugs de software lhes custam aproximadamente 60 bilhões de dólares por ano \cite{aniche_2014}. Sendo assim, essa etapa além de aumentar a confiabilidade e manutenção do sistema, evita que bugs sejam inseridos.

O desenvolvimento deste projeto envolve o uso de TDD, que se mostra essencial para o desempenho no desenvolvimento. Além disso segundo \citeonline{aniche_2014}, o TDD assegura a implementação dos requisitos funcionais elicitados, pois sua abordagem promove a utilização das especificações dos requisitos diretamente nos testes.

A estratégia de testes segue a tradicional pirâmide de testes, conforme descrito por \citeonline{cohn_2010}. Nesse modelo, a base é composta por testes unitários, que são executados de forma automatizada para garantir a estabilidade e a correção de componentes individuais do sistema. Para o backend em Java, optou-se pelo JUnit 5, enquanto para o frontend em JavaScript, optou-se pelo Jest. Esses frameworks são escolhidos pela facilidade de implementação de mocks e pelo suporte ao snapshot testing. Além disso, a abordagem inclui testes de caixa branca, testes de integração e testes de sistema para garantir a cobertura e a qualidade em diferentes níveis da aplicação.

Embora o software teóricamente funcione em multiplos sistemas operacionais é esperado que este projeto não irá testar todos os dispositivos. Sendo assim, para realização de testes de usabilidade serão utilizados iPhones e Androids porém com número limitado de aparelhos.

% \subsection{Plano de Teste}

% Para garantir a cobertura dos requisitos elicitados no projeto proposto, os seguintes tópicos descrevem a configuração dos testes realizados com base nos requisitos funcionais identificados.

